% !TeX root = ../papers/nc-p3.tex
\section{第三の疎長方形と三主成分の正則開集合}
\label{sec:parallel-third-rectangle}
\begin{theorem}\label{thm:parallel-third-rectangle}
前節の二形に含まれない未決定の台88--90は一つの有理三変数表示を持つ。
\[
 b\ne0,1,\quad AD\ne0,1,\quad L=AD-1,\quad M=(b-1)L,\quad C=1-(b-1)AD/b
\]
とする。薄い部分の非零成分は全て1とし、厚い部分の16成分を前節の $x_0,\ldots,x_{15}$ の順序で
\[
 \left(A,b,1,1,1,1,C,D,-1/L,-1/L,
 (AD-b)/(bL),D/L,A/L,-b/M,-1/L,-1/L\right)
\]
と取ればよい。指定した文字座標の対角倍率を除いてこの表示は一意である。
この族の一般点は標準左右AS正則、反対階数16、最小周期4で、無印の像の次元は3である。
さらにこの三変数族および前節の二つの主成分のそれぞれに、全閉点が標準ASとなる非空Zariski開集合がある。
そこでは各始点で四中心四次式上の自由階数が256である。
\end{theorem}
\begin{proof}[第一切断の特異な交差ブロック]
薄い四長方形の積条件は連結な一つのトーラスになり、残る係数は全て位置別対角作用で1にできる。
第一切断の厚い行列は
\[
 \begin{pmatrix}x_0&x_1&0&0\\x_6&x_7&1&1\\1&0&x_8&x_9\\1&0&x_{14}&x_{15}\end{pmatrix},
 \qquad
 \begin{pmatrix}x_2&x_3&1&0\\x_4&x_5&1&0\\1&1&x_{10}&x_{11}\\0&0&x_{12}&x_{13}\end{pmatrix}.
\]
第一行列の左上を $P$、第二行列の右下を $Q$ とする。全三台で $x_1 x_{13}\ne0$ である。
各々の指定ブロックを含む三次小行列式から $x_8\det P=-x_1$, $x_2\det Q=x_{13}$ となり両ブロックが可逆である。
従って他方の対角ブロックの四成分は、それぞれ一定値 $u=-x_1/\det P$, $v=x_{13}/\det Q$ である。
薄い部分の安定群は左上の八座標を同じ倍率、右下をその逆で変える。
実際その指数行列の非零Smith因子は $1^{11},2$ で、最後の符号は16座標に自明に作用する。
よって $v=1$ にできる。

$P=(A,b,C,D)$ と置く。第二切断は二つの交差ブロックが両色で可逆なので、左上
$(A,1,1,D)$ と $(1,b,C,1)$ も可逆である。その行列式を $L,M$ と書く。
前者のSchur補式は $u=-1/L$, $x_{12}=A/L$, $x_{11}=D/L$ を与える。
後者は $u=(1-b)/M$, $x_{13}=-b/M$, $x_{10}=(2-b-C)/M$ を与える。
従って $M=(b-1)L$, $C=1-(b-1)AD/b$。逆に両方の第一切断の行列式条件もこれで成立する。
両方の組に渡る残りの零成分候補は $C$ と $x_{10}$ である。
それぞれ $AD=b/(b-1)$ と $AD=b$ で消え、両方なら $b=AD=2$ となる。
どちらか一方だけが零の場合の二つの表示は台表の許した回転・反転・文字変更で同じ代表89となる。
$AD=0$ の境界の二台86,87は既に全係数で排除済みである。
\path{computations/verify_parallel_third_rectangle.py} が全係数表示と残余倍率を検算する。
\end{proof}

\begin{lemma}[中心四次式による有限長の自由性証明]\label{lem:parallel-residue-border}
体上の四頂点二次経路代数を $B$ とする。四つの中心四次式を割った商が、各始点で長さ0から13まで
$(1+z+z^2+z^3)^4$ の係数を持つとする。元の各始点の商次元が長さ14まで $\binom{n+3}{3}$ なら、
各始点は四式の生成する多項式環上で階数256の自由加群である。
一文字階数4と隣接階数6の最上テンソルが与えられていれば、標準Koszul左右AS性が従う。
\end{lemma}
\begin{proof}
中心商は長さ13で零だから、その後も零である。各始点で商の斉次基底の語を長さ12まで持ち上げ、
これらを形式的な基底とする $P=k[z_1,z_2,z_3,z_4]$ 上の自由加群 $F$ を作る（各変数の重み4）。
中心性と次数帰納で自然な写像 $F\to e_i B$ は全射である。次元仮定から長さ14までは同型である。
各基底に右から矢を掛けたものを、長さ13までのこの同型で $F$ に戻すことで、$P$ 線型な矢作用を定義する。
二次関係を基底に作用させると長さ14以下であり、$e_iB$ に写して零だから $F$ 上でも零となる。
従って右作用は $B$ を経由する。空語に基底語を作用させると長さ12以下で対応する形式基底が得られ、
空語に第 $j$ 中心式を作用させると $z_j$ 倍の空語となる。右作用の $P$ 線型性により任意の多項式倍の基底が空語から得られる。
よってこの右作用が先の全射の逆を与える。多項式変数も代数的に独立である。
最後は定理\ref{thm:klein-rational-period-four}の証明と同じである。二つの互いに素な中心変数により
単純から環への $\operatorname{Hom},\operatorname{Ext}^1$ が消え、両側の表示の双対完全性とEuler等式で中央項も完全となる。
\end{proof}

\begin{proof}[三主成分での開集合]
第三の長方形では $(A,D,b)=(1,2,2),(2,3,5)$ を取る。有理数体上で中心四次式を求め、素数1009で還元する。
両点で四中心式の商が四始点で補題の長さ13までの値を持ち、還元後の元の商も長さ14まで標準次元を持つ。
隣接関係の基底、四つの一文字縮約、最上テンソルはこの素数で所定の階数を持つ。
従ってまず剰余体そのもので補題が適用できる。

この証明から標数零へ移るとき、剰余体上の余核次元を標数零での下界として用いる必要はない。
$\mathbb Z_{(1009)}$ 上で単元小行列式により縮約空間の自由基底を選び、テンソルとその縮約による
四項の線型Koszul列を作る。連続微分の積は定義から零である。長さ $n-1$ までの経路加群が自由なら、
内部次数 $n$ のこの列で最後の商以外は自由である。剰余体上の完全性により、始めから単元小行列式で
分裂単射を順に除ける。最後の余核は表示により長さ $n$ の経路加群で、再び自由となる。
初期次数から帰納して全次数で標数零の正則性を得る。左右両向きで同じ議論ができる。
そこで有理中心四式の商の次元上界と正則Hilbert級数を比較すれば標数零でも各始点が自由である。
この二点を含む三変数族の幾何学的一般点は前節の次数別の開条件により正則である。
$(2,3,5)$ で反対切断は可逆、四位置の無限小作用の階数は61、
前節の部分trace二乗の跡は四始点で相異なり、無印次元と最小周期が従う。

前節の二つの主成分では、同節の $t=2,b=3,(A,B,C)=(7,1,5)$ の有理点でそれぞれ四次中心の文字別次元が
$(3,2,0,0)$、四式の商の長さ13までの次元上界が補題の値となる。こちらは既に一般点が正則なので、
半連続性による下界と中心による自由加群の上界から各有理点自体が中心四変数上自由で標準正則となる。

以上の四有理点それぞれで、頂点で零となる内部次数4の導分を六つの二次関係へのLeibniz則で調べると、
224列に対する階数下界は文字順に $188,184,192,192$ である。これは内導分と中心四次式のEuler倍でちょうど尽くされる値である。
定理\ref{thm:full-component-central-open}の持ち上げ論法をそれぞれの既約基底で適用する。
すなわち、正則点での完備局所族の次数別平坦性は上と同じ単元分裂による。
次数4の中心の基底を非零因子で取れば命題\ref{prop:euler-central-lifting}が適用できる。
完備化の忠実平坦性と五次までの定階数性から、中心の階数5の束を一つの開近傍に得る。
長さ13までの中心商の上界は有限個の開条件で保たれ、一般点のHilbert下界から開集合の全点で
四変数上の自由性と標準AS性が従う。これらの主成分の閉包の任意の標準点にも、
系\ref{cor:klein-pi-specializations}と同じ議論で標準恒等式 $s_{512}=0$ が特殊化する。
\path{computations/verify_parallel_residue_border.py} が商経路、中心、導分の計算を再現する。
\end{proof}
